
\Problem{LOL}

\Statement

En l'argot d'internet,
LOL, l'acrònim anglès de ``laughing out loud'' (rient sonorament)
s'utilitza freqüentment per descriure una situació suposadament divertida.
A veure si aquest problema us sembla divertit...

\medskip

Donats dos naturals $n$ i $m$,
cal omplir una matriu $n \times m$
amb els caràcters \q{L} i \q{O},
de tal manera que el nombre de \qq{LOL}s que contingui sigui el màxim possible,
comptant totes les aparicions horizontals, verticals i diagonals.

\medskip

Per exemple, per a $n = 3$ i $m = 7$
la solució òptima és
\begin{verbatim}
LLLLLLL
OOOOOOO
LLLLLLL
\end{verbatim}
amb 17 aparicions de \qq{LOL}.
Com un altre exemple, per a $n = 1$ i $m = 2$
\begin{verbatim}
OL
\end{verbatim}
és una de les quatre solucions possibles,
totes sense cap aparició de \qq{LOL}.


\Input

L'entrada consisteix en diversos casos,
cadascun amb $n$ i $m$.
Suposeu $1 \le n \le m$ i $n \cdot m \le 23$.


\Output

Per a cada cas,
escriviu el màxim nombre possible de \qq{LOL}s.


\Observation

No es valorarà cap solució que no sigui de força bruta.


\Sample
